% Argumenten las principales diferencias observadas en las dinámicas del pozo nito simétrico y no simétrico. Les recomiendo que varíen los parámetros y analicen diferentes situaciones que puedan ocurrir.

% ¿Cuáles conceptos aprendimos con este proyecto? ¿Cuáles habilidades desarrollamos con el desarrollo del proyecto?

Las principales diferencias es que en el simétrico, podemos aprovechar la simetría del problema y sigue la misma estructura para $|x|>a$. En cambio, para el pozo asimétrico, es un dolor de cabeza normalizarlo, no hay simetría, pero de cierta manera se parece mucho al pozo simétrico en algunso sentidos.

Para llegar a explorar el modelo terminando el planteamiento de la ecuación maestra, utilicé un Desmos para explorar como funcionaba y contar el número de estados atrapados al variar $c$ y $n$.


Sobre los conceptos que aprendimos fue realizar los videos, llegar a practicar nuevamente los grids para mostrar todos los eigenestados de una forma nítida. 

De la parte de desarollo numérico, pues es muy interesante y práctico incluir  variables adimensionales que sirven a resolver los problemas.  Esto mismo recuerdo que se aplica en termodinámica y simplifica mucho las cosas. Después de varios intentos, proponer $z$ fue muy gratificador.